\documentclass[twocolumn,preprintnumbers,amsmath,amssymb]{revtex4-2}
\usepackage[utf8]{inputenc}
\usepackage[T2A]{fontenc}
\usepackage{graphicx}
\usepackage{hyperref}
\usepackage{booktabs}
\usepackage{amsmath}
\usepackage{natbib}

\begin{document}

\title{AI-driven multi-proxy predictive system for solar activity over 55,000 years\\and probabilistic forecast of the next Grand Solar Minimum using Monte Carlo simulation}

\author{A.\,A.~Dmitriev}
\affiliation{Russian National Public Library for Science and Technology (GPNTB Russia), Moscow, Russia\\Researcher, Network Scientific Laboratory (NSL) for Digital Transformation and AI, GPNTB Russia}
\email{a@aadmitrieff.com}

\date{\today}

\begin{abstract}
We present a five-stage analysis pipeline (SOL-1) combining pre-telescopic sunspot records, auroral catalogs, eclipse corona morphology, radiocarbon ($^{14}$C, IntCal20), and cosmogenic beryllium ($^{10}$Be, GISP2) to identify and cross-validate solar activity periodicities over 55,000~years. Fourier analysis independently confirms the Gleissberg cycle ($\sim$88~yr) in 4 of 4 applicable proxies, the Suess/de~Vries cycle ($\sim$207~yr) in 3 of 3, and the Hallstatt cycle ($\sim$2300~yr) in 2 of 2 cosmogenic proxies. Cross-validation between $^{14}$C and $^{10}$Be yields Pearson $r = 0.47$. A six-component harmonic model, phase-fitted to five historically documented Grand Solar Minima (Oort through Dalton), is used to forecast future solar activity. Monte Carlo simulation (500 realizations with perturbed periods and amplitudes) yields a median next Grand Solar Minimum at $\sim$2126~CE (68\% CI: 2058--2206), with 63\% probability of a ``Grand'' event comparable to the Sp\"orer or Maunder Minima, and 95\% probability of at least a ``Deep'' minimum. Coupling with IPCC~AR6 anthropogenic warming scenarios (SSP1-2.6 through SSP5-8.5) indicates that the GSM will not offset global warming but will introduce additional climate instability through simultaneous cooling and warming stresses. We discuss implications for agriculture, energy systems, cosmic radiation exposure, and propose a distributed cluster-based preparedness framework.
\end{abstract}

\keywords{solar activity prediction, Monte Carlo simulation, spectral analysis, multi-proxy cross-validation, time series forecasting, computational astrophysics, Grand Solar Minimum, climate modeling}

\maketitle

% ================================================================
\section{Introduction}
\label{sec:intro}

Solar activity varies on multiple timescales, from the well-known $\sim$11-year Schwabe cycle to millennial-scale oscillations observable only through cosmogenic isotope proxies \citep{Usoskin2017}. Periods of anomalously low activity---Grand Solar Minima (GSMs)---have occurred repeatedly over the Holocene, with the Maunder Minimum (1645--1715~CE) being the most thoroughly documented \citep{Eddy1976}.

Forecasting the next GSM is of practical importance for climate adaptation planning, satellite operations, and agricultural policy. Previous forecasts range from the 2030s \citep{Zharkova2015,Shepherd2014} to the late 21st century \citep{Steinhilber2013}, depending on methodology. Physical dynamo models \citep{Zharkova2015} and empirical spectral models \citep{Steinhilber2013} have yielded divergent predictions, partly because they rely on different subsets of the observational record.

Here we present SOL-1, a systematic pipeline that integrates five independent solar activity proxies spanning 55,000~years, performs spectral analysis on each, cross-validates periodicities across proxies, and generates a probabilistic forecast of the next GSM using Monte Carlo methods.

% ================================================================
\section{Data Sources}
\label{sec:data}

\subsection{E1: Pre-telescopic sunspot records}

We use the catalog of \citet{Yau1988}, comprising 112~naked-eye sunspot observations from Chinese, Korean, and Japanese dynastic histories spanning 28~BCE to 1610~CE. Each record includes year, source chronicle, and qualitative description.

\subsection{E2: Auroral catalog}

The auroral database combines \citet{Silverman1992} and \citet{Keimatsu1970}, totaling approximately 2000~records from 567 to 1900~CE. Low-latitude aurorae ($< 45^\circ$) serve as proxies for intense geomagnetic storms driven by high solar activity.

\subsection{E3: Eclipse corona morphology}

We compiled 116~total solar eclipse records from \citet{Stephenson1997} and historical chronicles, classifying corona descriptions into three phases: \textit{minimum} (round, thin, symmetric), \textit{maximum} (extended streamers, complex structure), and \textit{transitional}. This provides direct observational evidence of the solar cycle phase at each eclipse.

\subsection{E4: Radiocarbon (IntCal20)}

The IntCal20 calibration curve \citep{Reimer2020} provides 9501~data points of atmospheric $\Delta^{14}$C over 55,000~years. $^{14}$C production is inversely proportional to solar activity via cosmic ray modulation by the heliospheric magnetic field.

\subsection{E5: Beryllium-10 (GISP2)}

The GISP2 ice core $^{10}$Be concentration record \citep{Finkel1997} provides 387~data points spanning 3--40~ka~BP. $^{10}$Be is an independent cosmogenic proxy governed by the same cosmic ray physics as $^{14}$C but with a different geochemical pathway.

% ================================================================
\section{Method}
\label{sec:method}

\subsection{Spectral analysis}

For each proxy, we performed the following steps:
\begin{enumerate}
    \item Resampling to a uniform time grid via linear interpolation
    \item Removal of long-term trends (geomagnetic, ocean reservoir) using a running mean filter (2000--3000~yr window)
    \item Application of a Hanning window to reduce spectral leakage
    \item Fast Fourier Transform (FFT) to obtain the power spectrum
    \item Peak detection using \texttt{scipy.signal.find\_peaks}
\end{enumerate}

\subsection{Cross-validation}

Periodicities were compared across all five proxies. A cycle is considered ``confirmed'' if its period appears (within 20\% tolerance) in multiple independent proxies. The $^{10}$Be--$^{14}$C correlation was computed on detrended, uniformly resampled data over the common time range.

\subsection{Forecast model}

Solar activity is modeled as a superposition of six harmonic components:
\begin{equation}
    S(t) = \sum_{i=1}^{6} A_i \cos\!\left(\frac{2\pi t}{P_i} + \phi_i\right)
\label{eq:model}
\end{equation}
where $P_i$ are the confirmed periods (11, 22, 88, 207, 1000, 2300~yr), $A_i$ are amplitude weights proportional to the number of cross-validations, and $\phi_i$ are phase offsets fitted to minimize $S(t)$ at the midpoints of five historically documented GSMs (Oort, Wolf, Sp\"orer, Maunder, Dalton).

Phase fitting uses Nelder-Mead optimization with 200 random restarts.

\subsection{Monte Carlo uncertainty estimation}

We perform 500 Monte Carlo realizations, perturbing:
\begin{itemize}
    \item Periods: $P_i \to P_i + \mathcal{N}(0, \sigma_{P_i})$, where $\sigma_{P_i}$ represents measurement uncertainty
    \item Amplitudes: $A_i \to A_i + \mathcal{N}(0, \sigma_{A_i})$
    \item Phases: $\phi_i \to \phi_i + \mathcal{N}(0, 0.15)$~rad
\end{itemize}

For each realization, we identify the deepest minimum in the interval 2026--2250~CE.

% ================================================================
\section{Results}
\label{sec:results}

\subsection{Confirmed periodicities}

Table~\ref{tab:cycles} summarizes the detected periodicities across all five proxies.

\begin{table}[h]
\caption{Solar activity periodicities detected in each proxy. Check marks indicate detection within 20\% of the nominal period.}
\label{tab:cycles}
\begin{ruledtabular}
\begin{tabular}{lccccc}
Cycle & E1 & E2 & E3 & E4 & E5 \\
\hline
Schwabe (11~yr) & \checkmark & \checkmark & --- & Nyq & Nyq \\
Gleissberg (88~yr) & \checkmark & \checkmark & \checkmark & \checkmark & Nyq \\
Suess (207~yr) & --- & \checkmark & \checkmark & \checkmark & Nyq \\
Eddy (1000~yr) & --- & --- & --- & \checkmark & \checkmark \\
Hallstatt (2300~yr) & --- & --- & --- & \checkmark & \checkmark \\
\end{tabular}
\end{ruledtabular}
\end{table}

The Gleissberg cycle is the most robustly confirmed, appearing independently in all four proxies with sufficient temporal resolution: E1 (88.2~yr), E2 (87.5~yr), E3 (88.8~yr), and E4 (87.6~yr).

\subsection{$^{10}$Be--$^{14}$C cross-validation}

Over the common time range (38,000--1,300~BCE), the detrended $^{10}$Be and $\Delta^{14}$C time series yield a Pearson correlation coefficient of $r = 0.47$, confirming that both proxies respond to the same cosmic ray flux modulated by solar activity.

\subsection{Grand Solar Minimum forecast}

The Monte Carlo ensemble yields:
\begin{itemize}
    \item Median next GSM: $\sim$2126~CE
    \item 68\% CI: 2058--2206~CE
    \item 90\% CI: 2030--2237~CE
    \item $P(\text{Grand Minimum}) = 63.4\%$
    \item $P(\text{Deep} + \text{Grand}) = 95.4\%$
    \item $P(\text{no significant minimum}) = 0.0\%$
\end{itemize}

The predicted median depth ($-0.545$ on our normalized activity index) corresponds to a GSM intermediate between the Wolf and Sp\"orer Minima in severity, with an estimated global cooling of $0.3$--$0.6^\circ$C.

\subsection{Validation on known GSMs}

The model correctly reproduces all five calibration GSMs as activity minima: Oort ($-0.170$), Wolf ($-0.355$), Sp\"orer ($-0.788$), Maunder ($-0.971$), and Dalton (marginal at $+0.040$). The Dalton Minimum, being the shallowest and shortest, is a known challenge for harmonic models.

\subsection{Independent GSM verification matrix (MAU-1)}

As a separate validation exercise, we checked each of the five known GSMs against all five SOL-1 proxies for direct evidence of reduced solar activity during the documented period. The verification matrix (Table~\ref{tab:mau1}) confirms that the Maunder Minimum is the most robustly supported (3/4 sources), while Be10 resolution is too coarse to resolve individual GSMs shorter than $\sim$200~yr. Overall, 10 of 18 testable cells are confirmed (56\%), consistent with the known heterogeneity of pre-telescopic records.

\begin{table}[h]
\caption{MAU-1 verification matrix: direct evidence of reduced solar activity during each GSM. $\checkmark$ = confirmed, $\times$ = not detected, $\cdot$ = no data in window.}
\label{tab:mau1}
\begin{ruledtabular}
\begin{tabular}{lcccccc}
GSM & E1 & E2 & E3 & E4 & E5 & Score \\
\hline
Oort & $\checkmark$ & $\times$ & $\cdot$ & $\checkmark$ & $\cdot$ & 2/3 \\
Wolf & $\checkmark$ & $\times$ & $\times$ & $\checkmark$ & $\cdot$ & 2/4 \\
Sp\"orer & $\times$ & $\times$ & $\cdot$ & $\checkmark$ & $\cdot$ & 1/3 \\
Maunder & $\checkmark$ & $\times$ & $\checkmark$ & $\checkmark$ & $\cdot$ & 3/4 \\
Dalton & $\checkmark$ & $\checkmark$ & $\times$ & $\times$ & $\cdot$ & 2/4 \\
\hline
\multicolumn{6}{l}{Total} & 10/18 \\
\end{tabular}
\end{ruledtabular}
\end{table}

% ================================================================
\section{Discussion}
\label{sec:discussion}

\subsection{Comparison with previous forecasts}

Our prediction ($\sim$2060--2200, 68\% CI) is most consistent with \citet{Steinhilber2013}, who obtained $2090 \pm 30$ using a similar cosmogenic isotope methodology. It is later and less severe than \citet{Zharkova2015}, whose two-component dynamo model predicted a Maunder-like minimum in the 2030s. We note that Solar Cycle~25 has exceeded Zharkova's predicted decline, which may indicate that the 2030s GSM onset is less likely than initially proposed.

\subsection{Interaction with anthropogenic warming}

Under the SSP2-4.5 scenario (median pathway, 45\% probability), global temperature is projected to reach $+2.7^\circ$C above preindustrial by 2100. A concurrent GSM would reduce this by $0.3$--$0.6^\circ$C, yielding a net warming of $+2.1$--$2.4^\circ$C. However, the primary impact would not be net cooling but rather \textit{increased climate variability}: the opposing forcings would amplify seasonal and regional temperature extremes.

\subsection{Practical implications}

We identify four domains of significant impact:
\begin{enumerate}
    \item \textbf{Agriculture:} Crop yield reductions of 15--30\% due to temperature instability and altered precipitation patterns
    \item \textbf{Energy:} 2--5\% reduction in photovoltaic output coinciding with increased heating demand
    \item \textbf{Space weather:} 15--30\% increase in galactic cosmic ray flux, affecting aviation, satellite electronics, and astronaut safety
    \item \textbf{Geopolitics:} Increased migration pressure and resource competition, particularly in equatorial regions experiencing simultaneous heat stress and food insecurity
\end{enumerate}

\subsection{Limitations}

\begin{enumerate}
    \item Solar cycles are quasi-periodic; a purely harmonic model does not capture nonlinear dynamo dynamics
    \item Phase coherence degrades over centuries, limiting prediction accuracy beyond $\sim$200~years
    \item Amplitude weights are assigned heuristically based on cross-validation counts
    \item The model does not account for stochastic ``grand minimum triggering'' mechanisms
\end{enumerate}

% ================================================================
\section{Conclusion}
\label{sec:conclusion}

The SOL-1 pipeline demonstrates that five independent solar activity proxies converge on a consistent set of periodicities (Gleissberg, Suess, Eddy, Hallstatt), enabling a data-driven forecast of the next Grand Solar Minimum. The probability of a significant minimum before 2200 exceeds 95\%, with a 63\% chance of a Grand event comparable to historical GSMs that produced measurable climate impacts.

The practical window for preparation is the next 25--50 years. We advocate for a distributed, cluster-based preparedness framework emphasizing local food and energy autonomy, international mutual-aid protocols, and adaptive infrastructure designed for climate instability rather than simple warming or cooling.

All code, data, and results are available at \url{https://github.com/abelbleoworld-blip/astro-dating} under CC~BY~4.0.

\begin{acknowledgments}
This work uses data from IntCal20 \citep{Reimer2020}, the GISP2 ice core program \citep{Finkel1997}, and historical catalogs compiled by \citet{Yau1988}, \citet{Silverman1992}, and \citet{Stephenson1997}. Computational analysis was performed using Python with NumPy, SciPy, and Matplotlib.
\end{acknowledgments}

\bibliography{sol1_refs}

\end{document}
